% Options for packages loaded elsewhere
% Options for packages loaded elsewhere
\PassOptionsToPackage{unicode}{hyperref}
\PassOptionsToPackage{hyphens}{url}
%
\documentclass[
  spanish,
  11pt,
  letterpaper,
  letterpaper,
  12pt]{book}
\usepackage{xcolor}
\usepackage[left=2.5cm,right=2cm,top=2cm,bottom=2cm]{geometry}
\usepackage{amsmath,amssymb}
\setcounter{secnumdepth}{-\maxdimen} % remove section numbering
\usepackage{iftex}
\ifPDFTeX
  \usepackage[T1]{fontenc}
  \usepackage[utf8]{inputenc}
  \usepackage{textcomp} % provide euro and other symbols
\else % if luatex or xetex
  \usepackage{unicode-math} % this also loads fontspec
  \defaultfontfeatures{Scale=MatchLowercase}
  \defaultfontfeatures[\rmfamily]{Ligatures=TeX,Scale=1}
\fi
\usepackage{lmodern}
\ifPDFTeX\else
  % xetex/luatex font selection
\fi
% Use upquote if available, for straight quotes in verbatim environments
\IfFileExists{upquote.sty}{\usepackage{upquote}}{}
\IfFileExists{microtype.sty}{% use microtype if available
  \usepackage[]{microtype}
  \UseMicrotypeSet[protrusion]{basicmath} % disable protrusion for tt fonts
}{}
\makeatletter
\@ifundefined{KOMAClassName}{% if non-KOMA class
  \IfFileExists{parskip.sty}{%
    \usepackage{parskip}
  }{% else
    \setlength{\parindent}{0pt}
    \setlength{\parskip}{6pt plus 2pt minus 1pt}}
}{% if KOMA class
  \KOMAoptions{parskip=half}}
\makeatother
% Make \paragraph and \subparagraph free-standing
\makeatletter
\ifx\paragraph\undefined\else
  \let\oldparagraph\paragraph
  \renewcommand{\paragraph}{
    \@ifstar
      \xxxParagraphStar
      \xxxParagraphNoStar
  }
  \newcommand{\xxxParagraphStar}[1]{\oldparagraph*{#1}\mbox{}}
  \newcommand{\xxxParagraphNoStar}[1]{\oldparagraph{#1}\mbox{}}
\fi
\ifx\subparagraph\undefined\else
  \let\oldsubparagraph\subparagraph
  \renewcommand{\subparagraph}{
    \@ifstar
      \xxxSubParagraphStar
      \xxxSubParagraphNoStar
  }
  \newcommand{\xxxSubParagraphStar}[1]{\oldsubparagraph*{#1}\mbox{}}
  \newcommand{\xxxSubParagraphNoStar}[1]{\oldsubparagraph{#1}\mbox{}}
\fi
\makeatother


\usepackage{longtable,booktabs,array}
\usepackage{calc} % for calculating minipage widths
% Correct order of tables after \paragraph or \subparagraph
\usepackage{etoolbox}
\makeatletter
\patchcmd\longtable{\par}{\if@noskipsec\mbox{}\fi\par}{}{}
\makeatother
% Allow footnotes in longtable head/foot
\IfFileExists{footnotehyper.sty}{\usepackage{footnotehyper}}{\usepackage{footnote}}
\makesavenoteenv{longtable}
\usepackage{graphicx}
\makeatletter
\newsavebox\pandoc@box
\newcommand*\pandocbounded[1]{% scales image to fit in text height/width
  \sbox\pandoc@box{#1}%
  \Gscale@div\@tempa{\textheight}{\dimexpr\ht\pandoc@box+\dp\pandoc@box\relax}%
  \Gscale@div\@tempb{\linewidth}{\wd\pandoc@box}%
  \ifdim\@tempb\p@<\@tempa\p@\let\@tempa\@tempb\fi% select the smaller of both
  \ifdim\@tempa\p@<\p@\scalebox{\@tempa}{\usebox\pandoc@box}%
  \else\usebox{\pandoc@box}%
  \fi%
}
% Set default figure placement to htbp
\def\fps@figure{htbp}
\makeatother



\ifLuaTeX
\usepackage[bidi=basic]{babel}
\else
\usepackage[bidi=default]{babel}
\fi
% get rid of language-specific shorthands (see #6817):
\let\LanguageShortHands\languageshorthands
\def\languageshorthands#1{}


\setlength{\emergencystretch}{3em} % prevent overfull lines

\providecommand{\tightlist}{%
  \setlength{\itemsep}{0pt}\setlength{\parskip}{0pt}}



 


\usepackage{setspace}
\usepackage{titlesec}
\setcounter{secnumdepth}{0}
\titleformat{\chapter}[display]{\centering\normalfont\bfseries}{}{0pt}{}
\titleformat{\section}{\raggedright\normalfont\bfseries}{}{0pt}{}
\titleformat{\subsection}{\raggedright\normalfont\bfseries\itshape}{}{0pt}{}
\titleformat{\subsubsection}[runin]{\normalfont\bfseries}{}{0pt}{}[.]
\titlespacing{\subsubsection}{1.27cm}{*1}{1em}
\titleformat{\paragraph}[runin]{\normalfont\bfseries\itshape}{}{0pt}{}[.]
\titlespacing{\paragraph}{1.27cm}{*1}{1em}
\microtypesetup{expansion=false}
\doublespacing
\makeatletter
\@ifpackageloaded{caption}{}{\usepackage{caption}}
\AtBeginDocument{%
\ifdefined\contentsname
  \renewcommand*\contentsname{Tabla de contenidos}
\else
  \newcommand\contentsname{Tabla de contenidos}
\fi
\ifdefined\listfigurename
  \renewcommand*\listfigurename{Listado de Figuras}
\else
  \newcommand\listfigurename{Listado de Figuras}
\fi
\ifdefined\listtablename
  \renewcommand*\listtablename{Listado de Tablas}
\else
  \newcommand\listtablename{Listado de Tablas}
\fi
\ifdefined\figurename
  \renewcommand*\figurename{Figura}
\else
  \newcommand\figurename{Figura}
\fi
\ifdefined\tablename
  \renewcommand*\tablename{Tabla}
\else
  \newcommand\tablename{Tabla}
\fi
}
\@ifpackageloaded{float}{}{\usepackage{float}}
\floatstyle{ruled}
\@ifundefined{c@chapter}{\newfloat{codelisting}{h}{lop}}{\newfloat{codelisting}{h}{lop}[chapter]}
\floatname{codelisting}{Listado}
\newcommand*\listoflistings{\listof{codelisting}{Listado de Listados}}
\makeatother
\makeatletter
\makeatother
\makeatletter
\@ifpackageloaded{caption}{}{\usepackage{caption}}
\@ifpackageloaded{subcaption}{}{\usepackage{subcaption}}
\makeatother
\usepackage{bookmark}
\IfFileExists{xurl.sty}{\usepackage{xurl}}{} % add URL line breaks if available
\urlstyle{same}
\hypersetup{
  pdflang={es-CO},
  hidelinks,
  pdfcreator={LaTeX via pandoc}}


\author{}
\date{}
\begin{document}
\frontmatter


\mainmatter
\begin{titlepage}
    \thispagestyle{empty}
    \begin{center}
        {\huge \textbf{Universidad Nacional Mayor de San Marcos}} \\
        \vspace{3mm}
        {\large \textbf{Decana de América. Universidad del Perú}} \\

        \vspace{3mm}

        \includegraphics[width=.5\linewidth]{caratula/logo_unmsm.png}
        \vspace{3mm}

        {\large \textbf{Facultad de Ciencias Económicas}}\\
        \vspace{1mm}
        \large{Escuela Profesional de Economía Internacional}

        \vspace{2mm}
        \begin{spacing}{1.5}
            {\large \textsc{Evaluación del nivel de sincronización de los ciclos económicos en la región andina (Perú, Bolivia, Ecuador y Colombia) en el periodo de 2000 a 2019}}
        \end{spacing}

        \vspace{5mm}
        {\large \textbf{Bach. Nick Xavier Hurtado Valladolid}} \\
        \vspace{5mm}
        {\large Proyecto de Tesis para obtener el grado de Licenciado en Economía}
        \vfill
        {\huge \textbf{2026}}
    \end{center}
\end{titlepage}

\frontmatter

\chapter{Dedicatoria}\label{dedicatoria}

\emph{A mi familia, por el constante apoyo}

\emph{A mi futura novia, Ximena, por su inconmesurable paciencia en mis
noches de inmsonio para el desarrollo de esta tesis}

\chapter{Agradecimientos}\label{agradecimientos}

\chapter{Resumen}\label{resumen}

\chapter{Abstract}\label{abstract}

\tableofcontents
\listoftables
\listoffigures

\mainmatter

\chapter{Introducción}\label{introducciuxf3n}

Comprenderá la descripción panorámica de los capítulos, guardando una
secuencia lógica de tal manera que su lectura brinde una imagen
coherente y resumida del Diseño del Proyecto de investigación a
desarrollar, que además incluye los ANTECEDENTES informativos y/o
bibliográficos, también donde se indica la IMPORTANCIA de la
investigación a realizarse.

\chapter{Definición y Planteamiento del
Problema}\label{definiciuxf3n-y-planteamiento-del-problema}

\section{Identificación del
Problema}\label{identificaciuxf3n-del-problema}

\section{Delimitación del Problema}\label{delimitaciuxf3n-del-problema}

\section{Formulación del Problema}\label{formulaciuxf3n-del-problema}

\subsection{Problema General}\label{problema-general}

¿Cuál es el nivel de sincronización de los ciclos económicos de Perú,
Bolivia, Ecuador y Colombia durante el periodo 2000-2019?

\subsection{Problemas específicos}\label{problemas-especuxedficos}

\begin{itemize}
\item
  ¿Existe un factor cíclico común que explique de manera
  estadísticamente significativa una parte relevante de las
  fluctuaciones cíclicas del PBI de estos cuatro países?
\item
  ¿Cuál es el grado de sensibilidad de cada uno de los cuatro países
  respecto a dicho factor común?
\item
  ¿Qué proporción de la varianza cíclica de cada país se explica por el
  factor común, y qué proporción corresponde a factores idiosincráticos
  o propios de cada economía?
\end{itemize}

\section{Justificación del
Problema}\label{justificaciuxf3n-del-problema}

\subsection{Justificación Teórica}\label{justificaciuxf3n-teuxf3rica}

La presente investigación desea amplear el conocimiento en el área de la
integración económica a nivel regional. A pesar de distintos
planteamientos en la literatura, no se ha realizado los esfuerzos
suficientes en las investigaciones que relacionan la integración
económica con una sincronización ciclica. Por otro lado, la literatura
relacionada a la integracion económica buscan la aplicación de teorías a
paises latinoamericanos asociados en bloques como el Mercosur o en
comparaciones agregadas de la región, lo cual deja un vacio entorno a la
aplicación en la Comunidad Andina. Esta investigación se ubica en la
intersección de ambos vacíos, aportando evidencia específica sobre el
bloque andino mediante un enfoque de sincronización económica.

\subsection{Justificación Práctica}\label{justificaciuxf3n-pruxe1ctica}

Los resultados de esta investigación tiene gran utilidad para diferentes
actores económicos. En un primer lugar, para los actores de políticas
económicas, en tanto pueden conocer los efectos de coordinaciones
macroeconómicas y las posibles respuestas que pueden desembocar gracias
al conocimiento de el grado de interaccion entre los países de la
Comunidad Andina.

En segundo lugar, para los actores financieros de la región, quienes
pueden estimar riesgos provenientes de una relacion de actividades
económicas con las economias de la comunidad andina.

\subsection{Justificación
Metodológica}\label{justificaciuxf3n-metodoluxf3gica}

Se propone un diseño metodológico de dos etapas: En primer lugar, la
extracción del componente cíclico del producto bruto interno (PBI) de
cada país mediante el filtro de Hodrik-Prescott, de uso extendido por su
capacidad práctica. Como segundo paso, se utilizará el filtro de Kalman
para la estimación de un ciclo común regional.

\chapter{Hipótesis}\label{hipuxf3tesis}

\section{Hipotesis General}\label{hipotesis-general}

Estimar el nivel de sincronización de los ciclos económicos de Perú,
Bolivia, Ecuador y Colombia durante el periodo 2000-2019

\section{Hipotesis específica}\label{hipotesis-especuxedfica}

\begin{itemize}
\item
  Estimar un factor cíclico común regional que explique las
  fluctuaciones cíclicas del PBI de los cuatro países mediante un modelo
  de factor dinámico.
\item
  Determinar el grado de sensibilidad (carga factorial) de cada país
  respecto al factor cíclico común estimado.
\item
  Calcular la proporción de la varianza cíclica de cada país explicada
  por el factor común y por su componente idiosincrático.
\end{itemize}

\chapter{Objetivos}\label{objetivos}

\section{Objetivos General}\label{objetivos-general}

Estimar el grado de sincronización de los ciclos económicos de Bolivia,
Colombia, Ecuador y Perú mediante un modelo de espacio de estados con
filtro de Kalman, en el periodo {[}2000 - 2019{]}.

\section{Objetivos Específicos}\label{objetivos-especuxedficos}

\begin{enumerate}
\def\labelenumi{\arabic{enumi}.}
\item
  Extraer el componente cíclico del PBI real de cada país mediante un
  modelo estructural de series de tiempo (tendencia local lineal + ciclo
  estocástico).
\item
  Estimar un factor cíclico común regional y las cargas factoriales de
  cada país mediante un modelo de factor dinámico.
\item
  Comparar los resultados obtenidos con métodos tradicionales (filtro
  HP, correlaciones, índice de concordancia de Harding-Pagan) para
  evaluar la robustez de las conclusiones.
\end{enumerate}

\chapter{Marco Teorico}\label{marco-teorico}

Debe explicitar los aspectos conceptuales y referenciales que explican
en forma teórica, el fenómeno a investigar. Es importante, por que sin
su elaboración no es posible el Proyecto de Investigación.

\chapter{Metodología}\label{metodologuxeda}

\section{Consideraciones
Metodológicas}\label{consideraciones-metodoluxf3gicas}

Se incluirá en forma general, los diferentes métodos y/o técnicas a
implementarse en el desarrollo del proyecto. Los métodos y/o técnicas
específicas de recopilación de información, selección y análisis de
éstos, debiendo incluirse en las respectivas etapas o fases a considerar

\chapter{Esquema del proyecto de
investigación}\label{esquema-del-proyecto-de-investigaciuxf3n}

Escogido el problema apropiado, se establecerá en forma concreta los
aspectos del contenido del proyecto y /o Trabajo de Investigación a
desarrollarse sobre la base de un esquema capitular tentativo. La
estructura interna del trabajo de Investigación es la parte más
importante, cada uno de los

aspectos debe estar itemizado en capítulos y/o subcapítulos, guardando
una secuencia lógica con la hipótesis a demostrar de tal forma que cada
paso a seguir permita llegar a conclusiones.

\backmatter

\chapter{Bibliografía}\label{bibliografuxeda}

\chapter{Anexos}\label{anexos}


\backmatter


\end{document}
